\newpage

\section*{\Huge Tóm tắt}

Đề tài ``Phát triển ứng dụng di động phục vụ nhân sự Trường Đại học'' hướng đến xây dựng ứng dụng MyHCMUT Mobile nhằm hỗ trợ nhân sự Trường Đại học Bách khoa -- ĐHQG-HCM tiếp cận và thực hiện các nghiệp vụ quản lý, điều hành trên thiết bị di động. Thay vì xây dựng lại các hệ thống nghiệp vụ, ứng dụng được phát triển theo hướng tích hợp với các hệ thống hiện hữu của Nhà trường, qua đó bổ sung một kênh truy cập phù hợp hơn với thiết bị di động.

Trong phạm vi đề tài, nhóm tập trung phát triển các chức năng liên quan đến quản lý hồ sơ nhân sự, nghỉ phép, công tác, văn bản, nhiệm vụ và lịch làm việc. Ứng dụng khai thác dữ liệu và quy trình nghiệp vụ từ các hệ thống HRM, iOffice và dịch vụ xác thực hiện hữu thông qua các cơ chế tích hợp tương ứng. Bên cạnh các giao diện được xây dựng trực tiếp trên ứng dụng di động, một số quy trình hiện hữu trên Web được tiếp tục khai thác khi cần thiết.

MyHCMUT Mobile được phát triển bằng Flutter và tổ chức theo các mô-đun nghiệp vụ. Trong quá trình thực hiện, nhóm tiến hành phân tích yêu cầu, thiết kế ứng dụng, xây dựng các chức năng và tích hợp với các dịch vụ hiện hữu của Nhà trường. Một số luồng nghiệp vụ đại diện được kiểm thử trên thiết bị Android và đối chiếu trạng thái xử lý với các hệ thống liên quan.

Kết quả của đề tài là phiên bản ứng dụng MyHCMUT Mobile cung cấp các giao diện và luồng tương tác trên thiết bị di động cho những nhóm nghiệp vụ thuộc phạm vi đề tài. Kết quả kiểm thử cho thấy các luồng đã được kiểm chứng có thể khai thác dữ liệu và quy trình từ các hệ thống hiện hữu. Đồng thời, các trường hợp chưa được kiểm chứng đầy đủ và những giới hạn của phiên bản hiện tại được xác định làm cơ sở cho việc tiếp tục hoàn thiện và đánh giá ứng dụng trước khi xem xét triển khai ở phạm vi rộng hơn.
